% Precision-corrected insertion draft, 2026-11-06.
% Source: results/revision_c2/{candidate_metrics.csv,e2e_stable/e2e_metrics.csv,exports/}.
% Integrate with supplementary_methods.tex; historical tables use different metrics.
\subsection{Results of automatic calibration and coefficient transfer}

All five datasets completed the calibration protocol and full-test evaluation.
The study includes 1,120 native validation configurations, 20 directed
source--target transfers, 100 calibration-label subset runs, and 200 coefficient
perturbation runs. Frozen-checkpoint inference evaluates 60 dataset--edge-vector
combinations on validation and 55 dataset--variant combinations on test.
All inference uses the corrected FP32 graph-encoder/BF16 LLM policy, with
validation selections recomputed before test inference. Ten additional
test diagnostics hold the previously selected edge coefficients fixed at the
corrected precision; they do not participate in configuration selection.
The same 2,100 validation rows and 2,100 test rows are reused across their
respective variants; these are not independent new examples. Query-level caching
avoids repeated generation for identical model inputs.

\begin{table}[t]
\centering
\caption{Test candidate MRR@20. Original uses the native soft-support formula;
the other columns use identical shared feature definitions. LODO excludes target
validation labels from coefficient fitting and regularization selection.
These comparisons retain exactly the same top-20 candidate membership.}
\label{tab:c2-candidate-calibration}
\begin{tabular}{lrrrr}
\hline
Dataset & Original & Shared equal & Shared learned & Shared LODO \\
\hline
PrimeKG  & 0.125326 & 0.125188 & 0.125326 & 0.126104 \\
PharmKG  & 0.028159 & 0.031883 & 0.031876 & 0.031860 \\
Hetionet & 0.095439 & 0.101269 & 0.115284 & 0.113579 \\
DRKG     & 0.046556 & 0.061426 & 0.060496 & 0.058612 \\
repoDB   & 0.048107 & 0.057366 & 0.064902 & 0.064902 \\
\hline
\end{tabular}
\end{table}

Candidate-level results show that a common calibration feature space can retain
or improve the point estimates of the native soft-support formula. However,
the equal-weight control already improves several datasets, so comparisons
against the original formula combine feature changes with coefficient learning.
Against the feature-matched equal-weight control, learned calibration changes
MRR by 0.000138, $-0.000007$, 0.014015, $-0.000930$, and 0.007536 on PrimeKG,
PharmKG, Hetionet, DRKG, and repoDB, respectively. The exploratory paired
query-cluster 95\% bootstrap interval excludes zero only for repoDB among
these five comparisons: [0.000402, 0.015176]. Thus, the results do not establish
that learning uniformly outperforms a common fixed weighting.

The source-only LODO coefficients produce MRR point estimates above the native
formula on all five datasets, but they do not uniformly improve over shared
equal weights. In particular, DRKG shared equal weights yield 0.061426 versus
0.058612 for LODO. The LODO study concerns coefficient transfer within the
existing benchmark pipelines and should not be interpreted as zero-shot
completion or generalization to unseen diseases.

Learning curves also differ by dataset. On PharmKG, mean test MRR changes from
0.028975 at 10\% of calibration query groups to 0.031856 at 100\%, with standard
deviations 0.004282 and 0.000020 across subset/CV seeds, respectively. PrimeKG is
unchanged at 0.125326 over the tested fractions. Hetionet and repoDB do not show
monotonic improvement as the calibration fraction increases. These seed-based
standard deviations describe calibration variability, not confidence intervals
for performance on independent test sets.

\begin{table}[t]
\centering
\caption{Generated-answer exact accuracy (\%) in the candidate-order arm with
the full source graph and a fixed checkpoint within each dataset. Accuracy is
separate from candidate MRR and adjusted MRR.}
\label{tab:c2-generated-accuracy}
\begin{tabular}{lrrrrr}
\hline
Dataset & N & Original & Equal & Learned & LODO \\
\hline
PrimeKG & 500 & 0.0 & 0.0 & 0.0 & 0.0 \\
PharmKG & 500 & 0.0 & 0.2 & 0.2 & 0.2 \\
Hetionet & 100 & 3.0 & 1.0 & 1.0 & 2.0 \\
DRKG & 500 & 1.4 & 0.8 & 2.0 & 1.8 \\
repoDB & 500 & 0.2 & 0.8 & 0.0 & 0.0 \\
\hline
\end{tabular}
\end{table}

End-to-end outcomes are mixed (Table~\ref{tab:c2-generated-accuracy}).
Higher candidate MRR does not ensure higher generated-answer accuracy. For
example, repoDB learned calibration improves candidate MRR but yields no exact
answers, compared with four correct answers among 500 rows for shared equal
weights. The paired accuracy difference is $-0.8$ percentage points, with an
exploratory 95\% interval of [$-1.702$, $-0.190$] percentage points. On DRKG,
learned calibration yields 2.0\% exact accuracy versus 0.8\% for equal weights;
the corresponding interval for the 1.2-point difference is [0.000, 2.469]
percentage points. These findings limit claims about improvements in generation.

\begin{table}[t]
\centering
\caption{Test adjusted MRR@20 in the graph-selection arm. Candidate order and
the compressed-graph budget are fixed per query. Invalid outputs retain the
explicit virtual-rank-one convention described in the protocol; these values
must not be read as generated-answer accuracy.}
\label{tab:c2-edge-calibration}
\begin{tabular}{lrrrr}
\hline
Dataset & Fuzzy & Fixed & Auto & Edge LODO \\
\hline
PrimeKG & 0.074687 & 0.074687 & 0.074676 & 0.074676 \\
PharmKG & 0.020971 & 0.020956 & 0.021317 & 0.020969 \\
Hetionet & 0.100087 & 0.100087 & 0.100087 & 0.100087 \\
DRKG & 0.046285 & 0.046285 & 0.046285 & 0.046285 \\
repoDB & 0.037776 & 0.037795 & 0.038159 & 0.036653 \\
\hline
\end{tabular}
\end{table}

For edge calibration, the numbers of vectors tied at the best validation
accuracy and adjusted MRR are 6/12, 2/12, 12/12, 12/12, and 1/12 for PrimeKG,
PharmKG, Hetionet, DRKG, and repoDB, respectively. The prespecified tie rule
selects the earliest vector, rather than using test results to break ties.
The four compressed-graph variants in Table~\ref{tab:c2-edge-calibration} have
the same exact-answer accuracy on PrimeKG, PharmKG, Hetionet, and DRKG:
0.0\%, 0.2\%, 3.0\%, and 1.4\%, respectively. On repoDB, original fuzzy,
shared fixed, and shared auto yield 0.4\%, while the newly selected edge-LODO
vector yields 0.2\%. Thus automatic edge calibration does not establish a
uniform exact-answer advantage. Equal metrics
under the tested vectors do not identify a unique optimum or prove equivalence
for other graphs, budgets, or checkpoints.

\paragraph{Precision correction and validation reselection.}
The CPU candidate-calibration results remain unchanged. The repeated GPU
validation search changes the selected automatic trial on PharmKG (0 to 3) and
repoDB (0 to 2), while PrimeKG, Hetionet, and DRKG retain trials 1, 0, and 0.
Source-only edge-LODO selections change on all five datasets: 0 to 2, 7 to 2,
7 to 2, 7 to 2, and 1 to 3, respectively. Trial identifiers are zero-based and
refer to the unchanged predefined search list. In particular, the old claim
that all twelve vectors tie on PharmKG is superseded: only two tie for the best
validation score at corrected precision. Even when exact test accuracy is
unchanged, individual generated answers and adjusted ranks can change.

Across the 23,100 test row--variant evaluations, 679 normalized generated
answers change, and adjusted MRR changes in 15 of the 55 test summaries. These
counts include both precision changes and adaptive graph reselection; they are
not counts of independent new test examples. Exact accuracy changes in only
one summary, repoDB edge-LODO. Holding its old selected coefficients fixed at
stable precision retains 0.4\% accuracy, whereas the newly selected coefficients
yield 0.2\%; this difference is therefore attributable to reselection within
the stable-precision comparison.

The encoder audit covers 10,765 distinct query/candidate-order/edge-multiset
inputs from the legacy and corrected validation/test protocols. It finds 22
nonfinite FP16 calls, all associated with PrimeKG. In each full-source PrimeKG
test variant, nine rows have nonfinite FP16 encoder outputs. Casting finite
FP32 encoder outputs to FP16 remains nonfinite in six rows of each such
variant. All audited FP32 outputs and their BF16 casts are finite. The 22
distinct calls include repeated query cases with different candidate orders;
they must not be interpreted as 22 independent affected test examples.

The updated workbook provides every legacy--stable variant comparison and
separates changes at fixed graph inputs from changes following reselection.
The fixed-old-coefficient diagnostic further decomposes adaptive-arm metric
changes into a precision-policy term and a reselection term. Neither term
isolates graph arithmetic from LLM arithmetic. Archived FP16 predictions are
retained for audit and are not mixed with corrected outputs in the reported
tables.

All results above are retrospective analyses of existing benchmark splits.
Bootstrap intervals are exploratory and are not adjusted for multiple
comparisons. Full configuration traces, coefficients, query-level predictions,
feature ablations, transfer matrices, sensitivity analyses, and paired
comparisons are provided by the accompanying reproducibility scripts.
